% ---
% title: AI Document Library — User Guide (v1, original)
% date: 2026-04-11
% updated: 2026-06-28
% type: document
% vendor: claude
% model: claude-sonnet-4-6
% tags: [user-guide, superseded, v1-original, historical]
% related: [USER-GUIDE.tex]
% ---
% SUPERSEDED 2026-06-28. This is the original USER-GUIDE.tex, written
% 2026-04-11 for a clipboard-copy-paste workflow: pasting MASTER-PROMPT.md
% and the four RESUME files into a chat box by hand. That workflow no
% longer exists — RESUME, CHECKPOINT, NOTE, COMMIT, and NEW PROJECT are
% now all file-based operations run by the ai-library-ops skill, and
% CLAUDE.md auto-loads MASTER-PROMPT.md in Claude Code / cloud / iOS
% sessions. This document is more stale than ARCHITECTURE.tex was —
% it is frozen at roughly v01/v02 of the project. Archived as the
% historical record of the original copy-paste-era design. See
% USER-GUIDE.tex at the library root for the current guide. Not
% maintained going forward.

\documentclass[12pt,letterpaper]{article}

\usepackage[utf8]{inputenc}
\usepackage[T1]{fontenc}
\usepackage{geometry}
\usepackage{listings}
\usepackage{xcolor}
\usepackage{fancyhdr}
\usepackage{titlesec}
\usepackage{parskip}
\usepackage{booktabs}
\usepackage{array}
\usepackage{mdframed}
\usepackage{hyperref}
\usepackage{enumitem}


\usepackage{tocloft}
\usepackage{tabularx}
\usepackage{longtable}
\usepackage{multirow}

\geometry{
  top=1.25in,
  bottom=1.25in,
  left=1.25in,
  right=1.25in
}

\definecolor{codebg}{RGB}{245,245,243}
\definecolor{codeborder}{RGB}{200,198,190}
\definecolor{notebg}{RGB}{255,252,235}
\definecolor{noteborder}{RGB}{200,185,100}
\definecolor{warnbg}{RGB}{255,245,245}
\definecolor{warnborder}{RGB}{200,100,100}

\lstset{
  basicstyle=\ttfamily\small,
  backgroundcolor=\color{codebg},
  frame=single,
  framesep=6pt,
  rulecolor=\color{codeborder},
  breaklines=true,
  breakatwhitespace=false,
  keepspaces=true,
  columns=flexible,
  xleftmargin=6pt,
  xrightmargin=6pt,
  aboveskip=8pt,
  belowskip=8pt
}

\mdfdefinestyle{notebox}{
  backgroundcolor=notebg,
  linecolor=noteborder,
  linewidth=1pt,
  innerleftmargin=12pt,
  innerrightmargin=12pt,
  innertopmargin=10pt,
  innerbottommargin=10pt,
  skipabove=10pt,
  skipbelow=10pt
}

\mdfdefinestyle{warnbox}{
  backgroundcolor=warnbg,
  linecolor=warnborder,
  linewidth=1pt,
  innerleftmargin=12pt,
  innerrightmargin=12pt,
  innertopmargin=10pt,
  innerbottommargin=10pt,
  skipabove=10pt,
  skipbelow=10pt
}

\pagestyle{fancy}
\fancyhf{}
\fancyhead[L]{\small\textit{AI Library --- User Guide}}
\fancyhead[R]{\small\texttt{Layer 1 Foundation}}
\fancyfoot[C]{\small\thepage}
\renewcommand{\headrulewidth}{0.4pt}

\titleformat{\section}{\large\bfseries}{}{0em}{\thesection\quad}[\vspace{2pt}\hrule\vspace{4pt}]
\titleformat{\subsection}{\normalsize\bfseries}{\thesubsection\quad}{0em}{}
\titleformat{\subsubsection}{\normalsize\itshape}{\thesubsubsection\quad}{0em}{}

\setcounter{secnumdepth}{3}
\setcounter{tocdepth}{3}

\setlist[itemize]{topsep=4pt, itemsep=3pt, parsep=0pt, leftmargin=1.4em}
\setlist[enumerate]{topsep=4pt, itemsep=3pt, parsep=0pt, leftmargin=1.6em}

\hypersetup{
  colorlinks=true,
  linkcolor=black,
  urlcolor=black,
  bookmarksnumbered=true
}

\renewcommand{\cftsecleader}{\cftdotfill{\cftdotsep}}

% ---------------------------------------------------------------

\begin{document}

% Title page
\begin{titlepage}
\centering
\vspace*{1.5in}
{\Huge\bfseries AI Document Library}\\[12pt]
{\Large User Guide}\\[6pt]
{\large Layer 1 Foundation --- Plain Text System}\\[24pt]
\hrule height 0.8pt
\vspace{16pt}
{\normalsize
A complete system for storing, versioning, navigating, and resuming\\
AI-generated work across any model, any vendor, and any time horizon.\\
No application required. No platform dependency. No expiry.
}
\vspace{16pt}
\hrule height 0.4pt
\vspace{32pt}
{\small Version 1.0 \quad --- \quad \today}
\vfill
{\footnotesize
\texttt{AI-Library/USER-GUIDE.tex}\\
Plain text. Open format. Read in any decade.
}
\end{titlepage}

\newpage
\tableofcontents
\newpage

% ---------------------------------------------------------------
\section{Philosophy and Foundational Principles}

\subsection{The Core Problem This System Solves}

Every AI-assisted knowledge workflow faces the same structural failure: context
is lost the moment a session ends. The AI that helped you develop an argument,
refine a draft, or explore a research question retains nothing when the
conversation closes. The next session begins from zero.

The naive solution is to rely on platform memory features, project workspaces,
or cloud-based AI tools that promise persistent context. These solutions fail in
the medium to long term because they are owned by vendors. Platform memory
features change, get paywalled, or disappear in product updates. Workspaces
created on one platform cannot be transferred to another. The organisation of
your work is held hostage to a service you do not control.

This system solves the problem at the foundation. Every piece of context,
every version of an artifact, every decision made, and every instruction for
resuming a project is stored in plain text files that you own, organised in a
folder structure you control, in a format that any human or any AI can read
without any intermediary tool.

\subsection{The File is the Truth}

The governing principle of this system is: \textit{the file is the truth, not the
application.} All intelligence about a project lives inside the files themselves,
not in a database, not in an application's state, not in a vendor's memory system.

This principle has several practical consequences:

\begin{itemize}
  \item Moving a file preserves all of its context because the context is embedded
        in the file via frontmatter and structured sections.
  \item Sharing a project means sharing a folder. The recipient needs no account,
        no application, and no explanation beyond ``read THREAD.md first.''
  \item Switching AI vendors mid-project costs nothing. The new AI reads the
        instructions file and picks up where the previous AI left off.
  \item The library cannot be lost due to a service outage, account suspension,
        pricing change, or company acquisition.
\end{itemize}

\subsection{Why Plain Text Survives Generations}

Plain text files from the 1970s remain perfectly readable today on any modern
computer without conversion, emulation, or special software. No other digital
format can make this claim. PDF requires a PDF renderer. Word documents require
Microsoft Office or a compatible application. Proprietary database formats require
the originating software. Plain text requires only the ability to read characters.

Markdown extends plain text with a small set of conventions for formatting
(headings, lists, emphasis, code blocks) that are equally legible when the
formatting is not rendered. A Markdown file read as raw text in a terminal
is as comprehensible as one rendered in a Markdown editor. This is not true
of any other formatted document standard.

YAML frontmatter --- the metadata blocks used in this system --- is itself
plain text. It requires no parser to be human-readable. An AI model reading
a file with YAML frontmatter understands the metadata without being instructed
on how to parse it.

\subsection{The Checkpoint Discipline}

The second governing principle is: \textit{externalise memory before it is lost,
not after.}

AI context windows compress and lose fidelity as conversations grow longer.
The model you are working with in a long session has a degraded understanding
of decisions made early in the conversation compared to its understanding of
recent exchanges. This compression is not visible to the user. The model will
continue to produce outputs that appear coherent while quietly misremembering
earlier decisions.

The checkpoint ritual exists to counteract this. At regular intervals, you
ask the AI to produce a complete state dump of everything it knows, a full
copy of the current artifact, and a structured instruction set for resuming
the work in a new session. The checkpoint is called before compression becomes
severe, not after.

The three artifacts produced at each checkpoint correspond to three different
types of knowledge:

\begin{enumerate}
  \item The \textbf{artifact} captures what has been built.
  \item The \textbf{context} captures what has been established and decided.
  \item The \textbf{instructions} capture how to continue.
\end{enumerate}

No single artifact contains all three types of knowledge. All three must be
saved together at each checkpoint.

\subsection{Why the System is Generational}

A project folder created today under this system will be fully navigable by
a human or an AI in twenty years, fifty years, or a hundred years, subject
to one condition: the files are stored on a medium that can be read.

The folder structure is self-documenting. \texttt{THREAD.md} narrates the
arc of the work. Version-numbered files make the progression explicit.
The frontmatter on each file provides provenance. The resume block in each
instructions file provides a re-entry point.

A person who has never seen this system before can open any project folder,
read \texttt{THREAD.md}, and understand what was done, why, and where it
stands --- without reading any other file first.

An AI model that has never been introduced to this system can receive the
master prompt at the start of a session and be fully operational within
the library immediately.

This combination --- human navigability and AI operability, both from
plain text alone --- is what makes the system generational.

% ---------------------------------------------------------------
\section{System Architecture}

\subsection{The Three Layers}

The AI Library system is designed in three layers of increasing automation.
This user guide covers Layer 1 only. Layers 2 and 3 are described briefly
for context but are not covered in this document.

\begin{center}
\begin{tabular}{@{}lp{4cm}p{6.5cm}@{}}
\toprule
\textbf{Layer} & \textbf{Name} & \textbf{Description} \\
\midrule
1 & Foundation & Plain text files, folder structure, master prompt,
                 manual checkpoint ritual. Full control from the prompt.
                 No tooling required beyond a text editor. \\
\addlinespace
2 & Session & Claude Projects or equivalent platform workspace.
              Persona loaded automatically. Session memory handled by platform.
              Layer 1 files remain the permanent record. \\
\addlinespace
3 & Automation & MCP (Model Context Protocol) client reads the library folder
                 directly. Checkpoints scripted. Navigation automated.
                 Layer 1 files remain the truth. \\
\bottomrule
\end{tabular}
\end{center}

\begin{mdframed}[style=notebox]
\textbf{Note:} Layer 1 is the only mandatory layer. Layers 2 and 3 add
convenience but do not change the fundamental structure. A library built on
Layer 1 remains fully functional if Layers 2 and 3 are never implemented.
\end{mdframed}

\subsection{The Folder Hierarchy}

The complete folder structure for the AI Library is as follows:

\begin{lstlisting}
AI-Library/
|-- MAP.md
|-- MASTER-PROMPT.md
|-- USER-GUIDE.tex
|-- projects/
|   |-- [project-slug-a]/
|   |   |-- THREAD.md
|   |   |-- persona.md
|   |   |-- v01--artifact.tex
|   |   |-- v01--context.md
|   |   |-- v01--instructions.md
|   |   |-- v02--artifact.tex
|   |   |-- v02--context.md
|   |   `-- v02--instructions.md
|   `-- [project-slug-b]/
|       |-- THREAD.md
|       |-- persona.md
|       `-- v01--artifact.md
|-- docs/
|   `-- 2026-04-11--brief-title--claude.md
|-- research/
|   `-- 2026-03-20--topic-notes--gpt.md
|-- creative/
|   `-- 2026-02-14--story-draft--claude.md
|-- code/
|   `-- 2026-01-09--auth-module--claude.py
`-- inbox/
    `-- (unsorted items)
\end{lstlisting}

\subsubsection{Root Files}

Three files live at the library root and are never moved:

\begin{itemize}
  \item \texttt{MAP.md} --- the master index of all items in the library.
  \item \texttt{MASTER-PROMPT.md} --- the system prompt pasted at the start
        of every AI session.
  \item \texttt{USER-GUIDE.tex} --- this document.
\end{itemize}

\subsubsection{The projects/ Folder}

Each project is a subfolder with a short, descriptive slug as its name.
Project slugs use lowercase letters and hyphens only.
No spaces, no uppercase, no special characters.

Every project folder contains:

\begin{itemize}
  \item \texttt{THREAD.md} --- the project spine. Described in full in Section 4.
  \item \texttt{persona.md} --- the persona prompt used for this project.
  \item One set of three versioned files per checkpoint (artifact, context, instructions).
\end{itemize}

\subsubsection{The Standalone Subfolders}

Items that are not part of a multi-session project are stored in one of four
subfolders based on content type:

\begin{center}
\begin{tabular}{@{}ll@{}}
\toprule
\textbf{Folder} & \textbf{Content} \\
\midrule
\texttt{docs/} & Written output: reports, briefs, essays, letters, analyses \\
\texttt{research/} & Research notes, summaries, literature reviews, reference material \\
\texttt{creative/} & Stories, copy, ideas, conceptual drafts, brainstorm outputs \\
\texttt{code/} & Scripts, modules, technical specifications, configuration files \\
\bottomrule
\end{tabular}
\end{center}

\subsubsection{The inbox/ Folder}

The inbox is a staging area for items that have not yet been classified
and placed in their correct location. New material is dropped into the inbox
first without any renaming or reorganisation. Classification happens
as a separate deliberate step, not at the moment of capture.

\begin{mdframed}[style=warnbox]
\textbf{Important:} Do not sort or rename files at the moment of capture.
The inbox exists precisely to prevent hasty classification decisions.
Sort the inbox during a dedicated maintenance session, not mid-workflow.
\end{mdframed}

% ---------------------------------------------------------------
\section{File Conventions}

\subsection{Naming Conventions}

\subsubsection{Project Files}

Project files use a version prefix followed by a type identifier:

\begin{lstlisting}
v[NN]--[type].[ext]
\end{lstlisting}

Where:
\begin{itemize}
  \item \texttt{[NN]} is a two-digit zero-padded version number:
        \texttt{v01}, \texttt{v02}\ldots\texttt{v09}, \texttt{v10}, \texttt{v11}
  \item \texttt{[type]} is one of: \texttt{artifact}, \texttt{context},
        \texttt{instructions}
  \item \texttt{[ext]} is the appropriate file extension for the content type
\end{itemize}

The two-digit padding is mandatory. Operating systems and file managers that
sort filenames lexicographically will place \texttt{v10} between \texttt{v1}
and \texttt{v2} if single-digit numbers are used. Two-digit padding ensures
correct chronological sort order up to version 99.

Three files are produced at each checkpoint. All three carry the same version
number. Examples for version 3:

\begin{lstlisting}
v03--artifact.tex
v03--context.md
v03--instructions.md
\end{lstlisting}

\subsubsection{Artifact File Extensions}

The artifact extension matches the content format:

\begin{center}
\begin{tabular}{@{}ll@{}}
\toprule
\textbf{Extension} & \textbf{Use} \\
\midrule
\texttt{.tex} & LaTeX documents \\
\texttt{.md} & Markdown documents \\
\texttt{.py} & Python code \\
\texttt{.js} & JavaScript code \\
\texttt{.txt} & Plain text with no formatting \\
\texttt{.html} & HTML documents \\
\bottomrule
\end{tabular}
\end{center}

Context and instructions files are always \texttt{.md} regardless of the
artifact format.

\subsubsection{Standalone Files}

Standalone files (not part of a project) use:

\begin{lstlisting}
YYYY-MM-DD--short-descriptive-slug--vendor.[ext]
\end{lstlisting}

Where:
\begin{itemize}
  \item \texttt{YYYY-MM-DD} is the ISO 8601 date of creation
  \item \texttt{short-descriptive-slug} is 2--5 lowercase hyphen-separated words
  \item \texttt{vendor} is one of: \texttt{claude}, \texttt{gpt},
        \texttt{gemini}, \texttt{other}
\end{itemize}

Examples:
\begin{lstlisting}
2026-04-11--interest-rate-brief--claude.md
2026-03-28--landing-page-copy-v3--gpt.md
2026-02-15--rust-auth-module--claude.py
2026-01-09--novel-chapter-4--gemini.md
\end{lstlisting}

Date-prefixed filenames sort chronologically by default in any file manager.
The vendor suffix provides immediate provenance without opening the file.

\subsection{Frontmatter Schema}

Every standalone file begins with a YAML frontmatter block at the very first
line of the file. No content precedes it. The block is delimited by triple
hyphens.

\begin{lstlisting}
---
title: Full Human-Readable Title
date: YYYY-MM-DD
updated: YYYY-MM-DD
type: document
vendor: claude
model: claude-sonnet-4-6
tags: [strategy, q2, analysis]
related: [2026-03-20--prior-brief--claude.md]
---
\end{lstlisting}

\subsubsection{Field Definitions}

\begin{longtable}{@{}lllp{5.5cm}@{}}
\toprule
\textbf{Field} & \textbf{Type} & \textbf{Required} & \textbf{Description} \\
\midrule
\texttt{title} & String & Yes & Full human-readable title of the item \\
\texttt{date} & ISO date & Yes & Date of original creation \\
\texttt{updated} & ISO date & Yes & Date of most recent meaningful edit \\
\texttt{type} & Enum & Yes & Content type (see below) \\
\texttt{vendor} & Enum & Yes & AI vendor that produced this item \\
\texttt{model} & String & No & Specific model name and version \\
\texttt{tags} & List & No & Free list of keyword tags \\
\texttt{related} & List & No & Filenames of related items \\
\bottomrule
\end{longtable}

\subsubsection{Type Values}

\begin{center}
\begin{tabular}{@{}ll@{}}
\toprule
\textbf{Value} & \textbf{Meaning} \\
\midrule
\texttt{document} & Written output: report, brief, analysis, essay \\
\texttt{research} & Research notes, summaries, reference material \\
\texttt{creative} & Creative writing, copy, conceptual work \\
\texttt{code} & Code, scripts, technical specifications \\
\texttt{context} & A standalone context bundle for resuming work \\
\bottomrule
\end{tabular}
\end{center}

\subsubsection{Vendor Values}

\begin{center}
\begin{tabular}{@{}ll@{}}
\toprule
\textbf{Value} & \textbf{Meaning} \\
\midrule
\texttt{claude} & Anthropic Claude (any version) \\
\texttt{gpt} & OpenAI ChatGPT (any version) \\
\texttt{gemini} & Google Gemini (any version) \\
\texttt{other} & Any other AI model or mixed provenance \\
\bottomrule
\end{tabular}
\end{center}

Project checkpoint files (\texttt{v[NN]--artifact}, \texttt{v[NN]--context},
\texttt{v[NN]--instructions}) do not use frontmatter. Their metadata is stored
in \texttt{THREAD.md}, which covers the entire project.

% ---------------------------------------------------------------
\section{Project Structure in Detail}

\subsection{The Versioned Triplet}

Every checkpoint in a project produces exactly three files.
They are always produced together, always carry the same version number,
and are always stored in the same project folder.

\begin{center}
\begin{tabular}{@{}lp{4cm}p{5.5cm}@{}}
\toprule
\textbf{File} & \textbf{Cadence} & \textbf{Contents} \\
\midrule
\texttt{v[NN]--artifact} & Every checkpoint & The full output in its current state.
  LaTeX, Markdown, code, or plain text. Never summarised, never truncated. \\
\addlinespace
\texttt{v[NN]--context.md} & Every checkpoint & Complete state dump.
  Everything established, decided, ruled out, and left open.
  Dense and comprehensive. \\
\addlinespace
\texttt{v[NN]--instructions.md} & Every checkpoint & Structured resume prompt.
  Goal, background, artifact state, decisions, open questions, next task, persona.
  Functional: a new AI can operate from this file alone. \\
\bottomrule
\end{tabular}
\end{center}

The three files change at different rates within a project. The artifact changes
every session. The context changes every session. The instructions file typically
changes less frequently --- it may be unchanged between versions 2 and 3 if no
new structural decisions were made. However, a copy is always saved at every
checkpoint even if unchanged, because storage is inexpensive and a gap in the
version sequence creates ambiguity.

\subsection{THREAD.md}

\texttt{THREAD.md} is the spine of a project. It is the first file any person
or AI should read when approaching an existing project. It contains:

\begin{enumerate}
  \item A description of the project: what it is, what it produces, for whom.
  \item A reference to the persona used.
  \item A chronological checkpoint log: one entry per version, explaining what
        happened, what changed, and what remains open.
  \item A resume section: which version to use and how to start a new session.
\end{enumerate}

The checkpoint log is the most important section. Each entry answers:

\begin{itemize}
  \item What triggered this checkpoint?
  \item What is the current state of the artifact?
  \item What has the AI established in its context?
  \item Did the instructions file change? If so, what shifted?
  \item What decisions were made during this session?
  \item What questions remain open?
\end{itemize}

A person reading \texttt{THREAD.md} top to bottom should be able to understand
the full arc of thinking across all sessions without opening any version file.

\subsubsection{THREAD.md Template}

\begin{lstlisting}
# Thread: [Full Project Name]
> Started: YYYY-MM-DD | Latest checkpoint: v[NN] | Status: [in-progress / complete / paused]

## What this project is

[2-4 sentences. What is being produced? What question is being answered?
What does the completed artifact look like? Who is it for?]

## Persona

The persona used for this project is defined in `persona.md`.
[Note any changes to the persona across versions here.]

---

## Checkpoint log

### v01 -- YYYY-MM-DD
**Triggered by:** [e.g. end of first session, outline complete, context length]
**Artifact:** [one sentence on what the artifact contains at this point]
**Context:** [one sentence on what the AI had established]
**Instructions:** FIRST VERSION -- [one sentence on what the instructions tell a new AI]
**Key decisions made:**
- [decision and brief rationale]
**Still open:**
- [open question]

---

### v02 -- YYYY-MM-DD
**Triggered by:** [what caused this checkpoint]
**Artifact:** [what changed since v01]
**Context:** [what new ground was covered]
**Instructions:** UNCHANGED | CHANGED -- [if changed, what shifted]
**Key decisions made:**
- [decision]
**Still open:**
- [question]

---

## How to resume this project

1. Read this file top to bottom to reorient.
2. Open `v[NN]--instructions.md` (latest version).
3. Start a new chat. Paste `MASTER-PROMPT.md` first.
4. Type: RESUME
5. Paste the full contents of `v[NN]--instructions.md`.
6. Optionally paste `v[NN]--artifact.[ext]` for full context.
\end{lstlisting}

\subsection{persona.md}

The persona file contains the system prompt you use for this project.
It describes what kind of AI you need: domain expertise, tone, audience,
what to avoid, and what to treat as given.

The persona is stable across sessions. It changes only if the nature of
the work fundamentally shifts. A persona written for a legal research
project should not need to change between sessions one and ten.

\subsubsection{persona.md Template}

\begin{lstlisting}
# Persona: [Project Name]

You are [role and domain description].

You have deep expertise in [specific subject area].

You write for [audience description].

Your tone is [tone description].

You avoid [list of things to avoid].

You treat the following as established and do not re-litigate them:
- [established fact or decision]
- [established fact or decision]

You are rigorous about [specific standards: sources, evidence, logic, etc.].

When uncertain, you say so explicitly rather than proceeding with confidence.
\end{lstlisting}

% ---------------------------------------------------------------
\section{MAP.md}

\subsection{Purpose}

\texttt{MAP.md} is the master index of the entire library. It serves three
functions simultaneously:

\begin{enumerate}
  \item \textbf{Directory (Yahoo-style):} Browsable categories with titled,
        described entries. A human can scan the map and navigate to any item.
  \item \textbf{Search index (Google-style):} Each entry has a one-line
        summary. An AI can read the entire map and know which file to open
        without asking.
  \item \textbf{Wiki (Wikipedia-style):} The Threads section shows how
        ideas relate and evolved across multiple files and sessions.
\end{enumerate}

\subsection{MAP.md Structure}

\begin{lstlisting}
# AI Library -- Master Map
> Last updated: YYYY-MM-DD | Total items: [N]

[One paragraph: what this library contains, who maintains it, 
 how to navigate it.]

---

## Recent
[Table of 10 most recently updated items. Updated after every session.]

| Date | Title | Type | Vendor |
|------|-------|------|--------|
| YYYY-MM-DD | [Title](path) | document | claude |

---

## Topics

### [Topic Name]
- [Title](relative/path/to/file.md) `tag` `tag`
  YYYY-MM-DD . vendor . type
  One sentence: what this item is and why it matters.

---

## Projects
- [Project Name](projects/slug/THREAD.md) `tag`
  Started YYYY-MM-DD . latest v[NN] . [status] . vendor
  One sentence on where the project currently stands.

---

## Threads
[Each Thread entry shows how an idea evolved across files.]

### Thread: [Idea or Project Name]
One paragraph describing the arc of this thread.

v01 [Title](path) -- YYYY-MM-DD
  |
  v  led to
v02 [Title](path) -- YYYY-MM-DD
  |
  v  branched
v03a [Title](path) -- YYYY-MM-DD
v03b [Title](path) -- YYYY-MM-DD

---

## Tags Index
| Tag | Count | Representative Items |
|-----|-------|----------------------|
| strategy | 0 | -- |

---

## Vendors
| Vendor | Model | Items |
|--------|-------|-------|
| Claude | claude-sonnet-4-6 | 0 |
\end{lstlisting}

\subsection{MAP.md Maintenance}

Update \texttt{MAP.md} after every session:

\begin{itemize}
  \item Add the new or updated item to the correct Topic section.
  \item Update the Recent table (add the new row at the top, remove the oldest if over 10).
  \item Update the project entry if the status or latest checkpoint changed.
  \item Update or extend a Thread entry if new work connects to existing ideas.
  \item Update the Tags index counts.
\end{itemize}

Use the \texttt{ADD TO MAP} command in the master prompt to generate correctly
formatted entries. You supply the title, path, and one-line summary.
The AI produces the formatted entry.

% ---------------------------------------------------------------
\section{The Checkpoint Ritual}

\subsection{What It Is}

The checkpoint ritual is the mechanism by which you externalise the state of
a working session before context degrades or the session ends. It is the
most operationally critical practice in the entire system.

A checkpoint produces four outputs:

\begin{enumerate}
  \item The \textbf{artifact} --- the complete current output
  \item The \textbf{context} --- the complete state of established knowledge
  \item The \textbf{instructions} --- the complete resume prompt for a future session
  \item The \textbf{thread entry} --- the log entry for \texttt{THREAD.md}
\end{enumerate}

The master prompt specifies the exact format for each of these outputs.
When you type \texttt{CHECKPOINT v[NN]}, the AI produces all four in
sequence in labelled code blocks ready for extraction.

\subsection{When to Call a Checkpoint}

Call a checkpoint when any of the following conditions are met:

\begin{itemize}
  \item The conversation has been running for more than 45--60 minutes.
  \item A significant structural decision has been made that should be preserved.
  \item The artifact has reached a meaningful milestone (outline complete,
        section complete, full draft complete).
  \item You are about to close the session for any reason.
  \item You sense that the AI's responses are becoming less coherent or
        are inconsistently referencing earlier decisions. This is a signal
        that context compression is occurring.
  \item You are about to change direction significantly and want to preserve
        the current state before branching.
\end{itemize}

\begin{mdframed}[style=warnbox]
\textbf{Important:} Do not wait until context compression has already degraded
the session significantly before calling a checkpoint. Call it as a precaution
when you sense it might be needed, not as a rescue operation after quality
has dropped. The checkpoint is cheap. A degraded context is expensive to recover.
\end{mdframed}

\subsection{Step-by-Step Checkpoint Procedure}

\begin{enumerate}
  \item Determine the version number. It is the current highest version
        number in the project folder plus one. If the project folder is
        empty, the first checkpoint is \texttt{v01}.

  \item Type in the chat: \texttt{CHECKPOINT v[NN]}
        where \texttt{[NN]} is the version number you determined.

  \item The AI will produce four labelled code blocks in sequence.
        Do not interrupt the output. Wait for all four to complete.

  \item Copy the contents of the \textbf{ARTIFACT} block.
        Open a new file in your text editor.
        Paste the contents.
        Save as \texttt{v[NN]--artifact.[ext]} in the project folder.

  \item Copy the contents of the \textbf{CONTEXT} block.
        Open a new file.
        Paste the contents.
        Save as \texttt{v[NN]--context.md} in the project folder.

  \item Copy the contents of the \textbf{INSTRUCTIONS} block.
        Open a new file.
        Paste the contents.
        Save as \texttt{v[NN]--instructions.md} in the project folder.

  \item Copy the contents of the \textbf{THREAD ENTRY} block.
        Open \texttt{THREAD.md} for this project.
        Paste the block at the bottom of the Checkpoint Log section.
        Update the header line to reflect the new latest checkpoint.
        Save \texttt{THREAD.md}.

  \item Update \texttt{MAP.md}:
        Update the project's entry to reflect the new checkpoint number.
        Update the Recent table.
        Save \texttt{MAP.md}.

  \item The checkpoint is complete.
        You may continue the session, or close it.
\end{enumerate}

\subsection{Versioning Philosophy}

Versions in this system are checkpoints, not releases. They do not represent
finished states. They represent the state of the work at a moment in time,
preserved for the purpose of resumption and reference.

Do not wait until a section or draft is ``finished'' to call a checkpoint.
Call it when you have made progress worth preserving, regardless of whether
the work is complete. An artifact at \texttt{v04} is not four times better
than \texttt{v01}. It is four saved states, each recoverable.

The \texttt{THREAD.md} log is where you record what each version represents.
Two versions that are structurally identical except for a single key decision
are worth differentiating precisely because \texttt{THREAD.md} will explain
what changed and why.

% ---------------------------------------------------------------
\section{The Master Prompt}

\subsection{Purpose}

The master prompt is the control mechanism for the entire Layer 1 system.
Pasted at the start of any AI session with any model, it establishes:

\begin{itemize}
  \item The AI's operational role
  \item The complete library structure
  \item All naming and formatting conventions
  \item The checkpoint ritual and output formats
  \item All navigation commands
  \item All behavioural constraints
\end{itemize}

The AI operates within your system from the first message of the session.
You do not explain the system. You do not repeat instructions. You simply
paste the prompt and begin.

\subsection{Where to Store It}

\texttt{MASTER-PROMPT.md} lives in the root of \texttt{AI-Library/}.
It is the single file you copy from at the start of every session.

\texttt{MASTER-PROMPT.tex} is the compiled LaTeX version of the same content,
stored for archival and reference.

Do not modify the master prompt mid-project. If you need to change it,
increment a version number in the filename
(\texttt{MASTER-PROMPT-v2.md}) and note the change date.

\subsection{Command Reference}

\begin{longtable}{@{}lp{9cm}@{}}
\toprule
\textbf{Command} & \textbf{Effect} \\
\midrule
\texttt{CHECKPOINT v[NN]} &
Produces all four checkpoint blocks in sequence: artifact, context,
instructions, and thread entry. Each in its own labelled code block. \\
\addlinespace
\texttt{RESUME} &
Reads the instructions file pasted immediately after.
Reconstructs the project state. Confirms understanding in one paragraph.
Asks one clarifying question only if necessary. Waits for direction. \\
\addlinespace
\texttt{SHOW MAP} &
Reads any MAP.md content pasted. Produces a structured summary of
library state: item counts, active topics, in-progress projects,
most recent updates, and observed gaps or inconsistencies. \\
\addlinespace
\texttt{SHOW THREAD} &
Reads any THREAD.md content pasted. Narrates the intellectual arc
of the project in plain language. Identifies turning points,
key decisions, and unresolved questions. Produces a coherent paragraph,
not a list. \\
\addlinespace
\texttt{NEW PROJECT [name]} &
Produces a complete blank THREAD.md for the named project and a
structured prompt to help the user write the persona. \\
\addlinespace
\texttt{ADD TO MAP [title] [path] [summary]} &
Produces a correctly formatted MAP.md entry and a Recent table row. \\
\bottomrule
\end{longtable}

% ---------------------------------------------------------------
\section{Workflow Walkthroughs}

\subsection{Starting a Brand New Project}

\begin{enumerate}
  \item Create the project folder:
        \texttt{AI-Library/projects/[project-slug]/}

  \item Open a new AI chat. Paste \texttt{MASTER-PROMPT.md}.

  \item Type: \texttt{NEW PROJECT [your project name]}

  \item The AI produces a blank \texttt{THREAD.md} and a persona prompt.
        Copy and save \texttt{THREAD.md} to the project folder.

  \item Answer the persona prompt questions. The AI produces \texttt{persona.md}.
        Save it to the project folder.

  \item Begin the working session. Type your opening prompt.

  \item Work until a checkpoint is warranted.

  \item Run the checkpoint ritual. Save the three versioned files.
        Update \texttt{THREAD.md}. Update \texttt{MAP.md}.
\end{enumerate}

\subsection{Resuming an Existing Project}

\begin{enumerate}
  \item Open \texttt{THREAD.md} for the project. Read the checkpoint log
        to reorient. Note the latest version number and the open questions.

  \item Open a new AI chat. Paste \texttt{MASTER-PROMPT.md}.

  \item Type: \texttt{RESUME}

  \item Paste the full contents of \texttt{v[NN]--instructions.md}
        (the latest version).

  \item Optionally, after the instructions, type:
        ``Here is the current artifact:''
        and paste the full contents of \texttt{v[NN]--artifact.[ext]}.

  \item The AI confirms its understanding and waits.

  \item Direct the work. Continue until the next checkpoint.
\end{enumerate}

\subsection{Organising Existing Material}

If you have existing AI outputs, chat exports, LaTeX drafts, and instruction
sets that predate this system:

\begin{enumerate}
  \item Drop all existing files into \texttt{AI-Library/inbox/}
        without renaming them.

  \item Open a new AI chat. Paste \texttt{MASTER-PROMPT.md}.

  \item For each file in the inbox, paste its contents and type:
        ``Help me create the correct frontmatter and filename for this item.''

  \item The AI produces the frontmatter block and the correct filename.

  \item Save the file with the correct name to the correct subfolder.

  \item For existing multi-session projects, create a project folder,
        identify the version ordering of your existing artifacts, rename them
        using the \texttt{v[NN]--} convention, and write a \texttt{THREAD.md}
        that reconstructs the arc as best you can from what exists.

  \item Work through the inbox systematically. Do not rush.
        A partially sorted library is still usable. An unsorted library is not.
\end{enumerate}

\subsection{Sharing a Project}

\begin{enumerate}
  \item Navigate to \texttt{AI-Library/projects/[project-slug]/}.

  \item Compress the folder to a ZIP archive.

  \item Send the ZIP to the recipient.

  \item Tell the recipient: ``Unzip this folder. Read \texttt{THREAD.md} first.
        To resume the work with any AI, paste \texttt{MASTER-PROMPT.md} from
        the AI-Library root and then type RESUME followed by the contents of
        the latest instructions file.''

  \item The recipient needs no account, no application, and no onboarding
        beyond those two instructions.
\end{enumerate}

If the recipient does not have a \texttt{MASTER-PROMPT.md} file, include a
copy of it in the ZIP alongside the project folder.

% ---------------------------------------------------------------
\section{Storage and Durability}

\subsection{Where to Store the Library}

The AI Library folder is storage-agnostic. It works identically on:

\begin{itemize}
  \item A local folder on your Mac with no sync
  \item iCloud Drive
  \item Google Drive
  \item Dropbox
  \item Any other cloud storage service
  \item An external hard drive
  \item A USB drive
\end{itemize}

The recommendation is to use the cloud storage service you already use
daily for other files. Consistency of location reduces the chance of
forgetting to save new files in the right place.

\subsection{Backup Strategy}

No single storage location is sufficient for a permanent library.
The following backup strategy is recommended:

\begin{itemize}
  \item \textbf{Primary:} A cloud-synced folder (iCloud, Google Drive,
        or Dropbox). Provides automatic versioning and cross-device access.

  \item \textbf{Secondary:} A quarterly copy to an external hard drive.
        Provides protection against account loss, cloud service failure,
        and accidental deletion.

  \item \textbf{Optional:} A git repository initialised in the
        \texttt{AI-Library/} folder. Provides full version history of
        every file with commit messages. Run \texttt{git init} once and
        \texttt{git commit} after each session.
\end{itemize}

\begin{mdframed}[style=notebox]
\textbf{On git:} Initialising a git repository in \texttt{AI-Library/} costs
nothing and requires no remote hosting. Even a local-only git history gives
you the ability to recover any version of any file at any point in time.
This is distinct from the version numbering in the checkpoint system,
which tracks intellectual progress. Git tracks every save, including
intermediate states.
\end{mdframed}

\subsection{What Survives and What Does Not}

\begin{center}
\begin{tabular}{@{}lcc@{}}
\toprule
\textbf{Scenario} & \textbf{Library survives} & \textbf{What you lose} \\
\midrule
Cloud provider shuts down & Yes (if secondary backup exists) & Nothing \\
Hard drive failure & Yes (if cloud sync active) & Nothing \\
Account suspension & Yes (if local copy exists) & Nothing \\
AI vendor shuts down & Yes & The live session only \\
You switch AI vendors & Yes & Nothing \\
Operating system changes & Yes & Nothing \\
Format becomes obsolete & Yes (plain text never does) & Nothing \\
\bottomrule
\end{tabular}
\end{center}

% ---------------------------------------------------------------
\section{Troubleshooting and Edge Cases}

\subsection{The AI Does Not Follow the Checkpoint Format}

Different AI models comply with structured output instructions at different
levels of reliability. If the checkpoint output is malformed or incomplete:

\begin{enumerate}
  \item After the malformed output, type: ``Redo the [block name] block.
        Use the exact format specified in the instructions. Output only
        that block in a code block labelled [LABEL].''

  \item If the model consistently fails to comply, simplify by asking for
        each block separately: ``Print only the ARTIFACT block now.''
        Then: ``Print only the CONTEXT block now.'' And so on.

  \item If using a model that does not support structured output well,
        copy the output as produced and reformat it yourself. The content
        matters more than the format. The format can be corrected at save time.
\end{enumerate}

\subsection{You Missed a Checkpoint and Context Has Degraded}

If a session ran too long without a checkpoint and the model is producing
inconsistent or low-quality outputs:

\begin{enumerate}
  \item Do not call a checkpoint from the degraded session.
        The context and instructions produced from a degraded session
        will themselves be degraded.

  \item Close the session.

  \item Open the latest instructions file from the previous checkpoint.

  \item Start a new session. Paste the master prompt. Resume from the
        last good checkpoint.

  \item Accept the loss of work done after the last checkpoint.
        This is the cost of not checkpointing early enough.
        It is recoverable. Future sessions will checkpoint earlier.
\end{enumerate}

\subsection{You Have Multiple Conflicting Versions of the Same File}

If you saved multiple versions of an artifact file at the same version
number (e.g. two files both named \texttt{v03--artifact.tex}):

\begin{enumerate}
  \item Do not delete either file.

  \item Rename one: \texttt{v03a--artifact.tex} and \texttt{v03b--artifact.tex}.

  \item Open \texttt{THREAD.md} and add a note to the v03 checkpoint entry
        explaining that a conflict occurred and which version (a or b) was
        used as the basis for subsequent work.

  \item Continue from the version you consider more current or correct.
\end{enumerate}

\subsection{The instructions.md Did Not Meaningfully Change}

It is common for the instructions file to remain largely unchanged across
two or three consecutive checkpoints, particularly in the middle of a long
project where the goal and structure are stable.

In this case, still save a copy of the instructions at the new version number.
In the \texttt{THREAD.md} entry, write: \textbf{Instructions: UNCHANGED}.
Do not skip saving the copy. A missing version in the sequence creates ambiguity
about whether the instructions were stable or simply not saved.

\subsection{Sorting a Large Inbox}

If the inbox contains a large number of unsorted files, sort them in batches
rather than all at once. Work through them by project first: identify all files
that belong to the same project, move and rename them together, create the
project folder and \texttt{THREAD.md}, and then move to the next project.

Leave standalone files for last. Standalone files require only a frontmatter
block and a correct filename. Project files require the more complex work of
establishing version ordering and writing a retrospective \texttt{THREAD.md}.

% ---------------------------------------------------------------
\section{Quick Reference}

\subsection{File Naming Cheat Sheet}

\begin{lstlisting}
Project files:
  v01--artifact.tex
  v01--context.md
  v01--instructions.md

Standalone files:
  YYYY-MM-DD--slug--vendor.ext
  2026-04-11--analysis-brief--claude.md

Root files (never move these):
  MAP.md
  MASTER-PROMPT.md
  USER-GUIDE.tex

Per-project files (never move these out of the project folder):
  THREAD.md
  persona.md
\end{lstlisting}

\subsection{Session Start Checklist}

\begin{itemize}
  \item[$\square$] Open new chat
  \item[$\square$] Paste \texttt{MASTER-PROMPT.md}
  \item[$\square$] Type \texttt{RESUME} and paste latest \texttt{v[NN]--instructions.md}
        (existing project) OR type opening prompt (new topic)
  \item[$\square$] Optionally paste latest artifact for full context
\end{itemize}

\subsection{Checkpoint Checklist}

\begin{itemize}
  \item[$\square$] Determine version number (current max + 1)
  \item[$\square$] Type \texttt{CHECKPOINT v[NN]}
  \item[$\square$] Save ARTIFACT block as \texttt{v[NN]--artifact.[ext]}
  \item[$\square$] Save CONTEXT block as \texttt{v[NN]--context.md}
  \item[$\square$] Save INSTRUCTIONS block as \texttt{v[NN]--instructions.md}
  \item[$\square$] Paste THREAD ENTRY block into \texttt{THREAD.md}
  \item[$\square$] Update \texttt{THREAD.md} header (latest checkpoint: v[NN])
  \item[$\square$] Update \texttt{MAP.md} (project entry and Recent table)
\end{itemize}

\subsection{Command Quick Reference}

\begin{center}
\begin{tabular}{@{}ll@{}}
\toprule
\textbf{Type this} & \textbf{To do this} \\
\midrule
\texttt{CHECKPOINT v[NN]} & Save full state (all four blocks) \\
\texttt{RESUME} + paste instructions & Pick up an existing project \\
\texttt{SHOW MAP} + paste MAP.md & Get library overview \\
\texttt{SHOW THREAD} + paste THREAD.md & Get project arc narrative \\
\texttt{NEW PROJECT [name]} & Start a new project \\
\texttt{ADD TO MAP [title] [path] [summary]} & Get a MAP.md entry \\
\bottomrule
\end{tabular}
\end{center}

\vspace{16pt}
\hrule height 0.4pt
\vspace{12pt}

\begin{center}
{\small\texttt{AI-Library / USER-GUIDE.tex} \quad Version 1.0 \quad \today}\\[4pt]
{\small Plain text. No vendor. No application required. No expiry.}
\end{center}

\end{document}
